\documentclass[11pt,letterpaper]{article}
\usepackage[T1]{fontenc}
\usepackage[utf8]{inputenc}
\usepackage{amsmath,amsthm,mathtools}
\usepackage{newpxtext,newpxmath}
\usepackage[margin=1in,headheight=15pt]{geometry}
\usepackage{microtype}
\usepackage{booktabs,array,longtable}
\usepackage[dvipsnames]{xcolor}
\usepackage{enumitem}
\usepackage{listings}
\usepackage{fancyhdr}
\usepackage{hyperref}
\usepackage[nameinlink,noabbrev]{cleveref}
\definecolor{Ink}{HTML}{20374B}
\definecolor{Link}{HTML}{235D77}
\definecolor{Soft}{HTML}{F3F5F7}
\hypersetup{colorlinks=true,linkcolor=Link,urlcolor=Link,citecolor=Link,
 pdftitle={A proof of the generating function for OEIS A199085},
 pdfauthor={OpenAI},pdfsubject={Third derivatives of parenthesized self-exponentials}}
\pagestyle{fancy}
\fancyhf{}
\fancyhead[L]{\small\textsc{Parenthesized self-exponentials}}
\fancyhead[R]{\small\textsc{OEIS A199085}}
\fancyfoot[C]{\thepage}
\renewcommand{\headrulewidth}{0.3pt}
\setlength{\parskip}{0.38em}
\setlength{\parindent}{1.1em}
\setlength{\emergencystretch}{2em}
\allowdisplaybreaks[1]
\numberwithin{equation}{section}
\newtheorem{theorem}{Theorem}[section]
\newtheorem{lemma}[theorem]{Lemma}
\newtheorem{proposition}[theorem]{Proposition}
\newtheorem{corollary}[theorem]{Corollary}
\theoremstyle{definition}
\newtheorem{definition}[theorem]{Definition}
\newtheorem{example}[theorem]{Example}
\theoremstyle{remark}
\newtheorem{remark}[theorem]{Remark}
\newcommand{\E}{\mathcal E}
\newcommand{\Sset}{\mathcal S}
\newcommand{\Dset}{\mathcal D}
\newcommand{\NN}{\mathbb N}
\newcommand{\ZZ}{\mathbb Z}
\newcommand{\pos}[1]{\left(#1\right)_{+}}
\newcommand{\GF}{A(z)}
\DeclareMathOperator{\depth}{depth}
\lstset{basicstyle=\small\ttfamily,breaklines=true,columns=fullflexible,
 keepspaces=true,frame=single,rulecolor=\color{black!20},backgroundcolor=\color{Soft},
 showstringspaces=false,tabsize=4,aboveskip=0.8em,belowskip=0.8em}

\title{\color{Ink}\LARGE\bfseries A proof of the generating function\\[0.2em]
for OEIS A199085\\[0.55em]
\large\normalfont Third derivatives of parenthesized self-exponentials}
\author{A self-contained mathematical article and computational companion}
\date{September 20, 2026}

\begin{document}
\maketitle
\begin{abstract}
Let $a(n)$ count the distinct third derivatives at $x=1$ of fully
parenthesized exponentiation expressions containing $n$ copies of $x$.
We prove
\[
 a(1)=a(2)=1,\qquad
 a(n)=\left\lfloor\frac{n^2-2}{3}\right\rfloor\quad(n\ge3),
\]
and hence the rational generating function conjectured in OEIS A199085.
The proof does not extrapolate a recurrence from computed terms.
Instead, it classifies all possible pairs of second and third derivatives.
At fixed leftmost nesting depth $k$, the normalized third derivatives
form the arithmetic progression
$k^2+2,\,k^2+4,\ldots,k^2+2(n-k-1)$, apart from one extremal singleton.
Only progressions of the same parity can overlap, and their ordered
endpoints permit exact counting without an inclusion--exclusion
expansion. A constructive realization proves that no values in these
progressions are missing. We also give recurrences, extremal values,
worked examples, and independently implemented exact-arithmetic checks.
\end{abstract}

\begingroup
\small
\setlength{\parskip}{0pt}
\setcounter{tocdepth}{1}
\makeatletter
\renewcommand{\l@section}{\@dottedtocline{1}{0em}{2.4em}}
\makeatother
\tableofcontents
\endgroup
\clearpage

\section{The problem and the result}
\label{sec:intro}
A fully parenthesized exponentiation expression is defined recursively:
$x$ is an expression, and if $G$ and $H$ are expressions, then $G^H$
is an expression. Its size is the number of occurrences of $x$.
Let $\E_n$ denote the expressions of size $n$. They are interpreted as
positive real-analytic functions on $(0,\infty)$, with
\[
 (G^H)(x)=\exp\bigl(H(x)\log G(x)\bigr).
\]
Thus the derivative at $1$ has no branch ambiguity. We count values,
not expressions: two different parenthesizations that give the same
third derivative contribute only one value.

Define
\begin{equation}\label{eq:defa}
 \Dset_n=\{F'''(1):F\in\E_n\},\qquad a(n)=|\Dset_n|,
 \qquad \GF=\sum_{n\ge1}a(n)z^n.
\end{equation}
The generating-function variable is $z$, to distinguish it from the
variable $x$ in the functions being differentiated.

The OEIS entry attributes the sequence to Vladimir Reshetnikov and the
conjectured generating function to Alois P. Heinz, dated November 2,
2011. A program credited to Vladimir Joseph Stephan Orlovsky, dated
February 1, 2012, evaluates the quadratic formula proved below
\cite{OEIS199085}. The entry was consulted on September 20, 2026;
its generating function was still labeled conjectural there.
This article supplies a proof from the defining expressions, without
making a separate claim of bibliographic priority.

\begin{theorem}[The generating function]\label{thm:main}
For the sequence in \eqref{eq:defa},
\begin{equation}\label{eq:closed}
 a(n)=
 \begin{cases}
  1,&n=1,2,\\[2pt]
  \displaystyle\left\lfloor\frac{n^2-2}{3}\right\rfloor,&n\ge3.
 \end{cases}
\end{equation}
Its ordinary generating function is
\begin{align}
 \GF
 &=\frac{z}{1-z}+\frac{z^3}{(1-z)^3}
       -\frac{z^7}{(1-z)^2(1-z^3)}\label{eq:gfdecomp}\\[2pt]
 &=\frac{z(1-z+z^2+2z^4-z^6)}{(1-z)^2(1-z^3)}\label{eq:gfpositive}\\[2pt]
 &=\frac{z(z^6-2z^4-z^2+z-1)}
 {z^5-2z^4+z^3-z^2+2z-1}.\label{eq:gfoeis}
\end{align}
The last expression is exactly the proposed OEIS generating function,
with its variable renamed to $z$.
\end{theorem}

The substantive step is not the rational-function simplification.
It is determining the complete set of derivative values and then
counting its overlaps. We prove a stronger classification in
\cref{thm:states}; \cref{sec:gf,sec:closed} then prove
\cref{thm:main}.

\section{The analytic reduction to two integer statistics}
\label{sec:jets}
Every expression $F$ satisfies
\begin{equation}\label{eq:firstjets}
 F(1)=1,\qquad F'(1)=1.
\end{equation}
Indeed these identities hold for $x$. If they hold for $G$ and $H$,
then $G^H(1)=1$ and
\[
 (G^H)'=G^H\left(H'\log G+H\frac{G'}G\right),
\]
which gives $(G^H)'(1)=1$. Structural induction proves
\eqref{eq:firstjets}.

\begin{lemma}[Derivative composition law]\label[lemma]{lem:composition}
Write
\[
 p(F)=F''(1),\qquad q(F)=F'''(1).
\]
For $F=G^H$,
\begin{align}
 p(F)&=p(G)+2,\label{eq:p}\\
 q(F)&=q(G)+3p(G)+3p(H)+3.\label{eq:q}
\end{align}
\end{lemma}
\begin{proof}
Set $L=\log G$ and $V=HL$, so that $F=e^V$.
Using \eqref{eq:firstjets}, at $x=1$ we obtain
\[
 L=0,\quad L'=1,\quad L''=p(G)-1,\quad
 L'''=q(G)-3p(G)+2.
\]
Consequently,
\begin{align*}
 V(1)&=0,\qquad V'(1)=1,\\
 V''(1)&=2H'(1)L'(1)+H(1)L''(1)=p(G)+1,\\
 V'''(1)&=3p(H)+3\bigl(p(G)-1\bigr)
              +q(G)-3p(G)+2\\
         &=q(G)+3p(H)-1.
\end{align*}
Finally,
\[
 (e^V)''=e^V\bigl(V''+(V')^2\bigr),\qquad
 (e^V)'''=e^V\bigl(V'''+3V'V''+(V')^3\bigr),
\]
which yield \eqref{eq:p} and \eqref{eq:q}.
\end{proof}

Introduce the normalized statistics
\begin{equation}\label{eq:normalized}
 k(F)=\frac{F''(1)}2,\qquad b(F)=\frac{F'''(1)}3.
\end{equation}
Since $k(x)=b(x)=0$, the composition law becomes
\begin{equation}\label{eq:state-law}
 \boxed{\quad k(G^H)=k(G)+1,\qquad
 b(G^H)=b(G)+2k(G)+2k(H)+1.\quad}
\end{equation}
It follows by induction that $k(F)$ and $b(F)$ are nonnegative integers.
In particular, every third derivative in question is divisible by $3$.
Multiplication by $3$ is injective, so counting the normalized values
$b(F)$ is equivalent to counting the actual derivatives.

A useful asymmetry is already visible: $b(H)$ does not occur in
\eqref{eq:state-law}. The exponent contributes its second derivative,
but not its third derivative. This is the analytic reason why two
statistics suffice.

\section{The left spine and its leaf budget}
\label{sec:spine}
View an expression as a rooted ordered binary tree: an internal node
represents exponentiation, its left subtree is the base, and its right
subtree is the exponent. The \emph{left spine} is the path obtained by
repeatedly following the left child from the root to a leaf.
Its length is the number of internal nodes on that path.

The first identity in \eqref{eq:state-law} says precisely that $k(F)$
is the length of the left spine. Equivalently, every non-leaf expression
has a unique decomposition
\begin{equation}\label{eq:spine}
 F=\bigl(\cdots((x^{H_1})^{H_2})\cdots\bigr)^{H_k},
 \qquad k=k(F),
\end{equation}
where $H_1$ is the exponent nearest the bottom of the spine.
If $m_i$ is the size of $H_i$, then
\begin{equation}\label{eq:leafbudget}
 n=1+\sum_{i=1}^k m_i,\qquad m_i\ge1.
\end{equation}
The single extra leaf is the bottom leftmost $x$.

\begin{lemma}[Spine formula]\label[lemma]{lem:spine}
For \eqref{eq:spine}, let
$s=\sum_{i=1}^k k(H_i)$. Then
\begin{equation}\label{eq:spineformula}
 F''(1)=2k,\qquad F'''(1)=3k^2+6s.
\end{equation}
\end{lemma}
\begin{proof}
After $i$ nodes of the spine have been assembled, the left-depth
statistic is $i$. The $i$th use of \eqref{eq:state-law} therefore adds
$2(i-1)+2k(H_i)+1$ to $b$. Starting from $b(x)=0$ gives
\[
 b(F)=\sum_{i=1}^k(2i-1)+2\sum_{i=1}^k k(H_i)
     =k^2+2s.
\]
The definition of $k(F)$ gives the second derivative as well.
\end{proof}

Each $H_i$ of size $m_i$ has at most $m_i-1$ internal nodes, and hence
$0\le k(H_i)\le m_i-1$. Moreover, $k(H_i)=0$ holds exactly when
$H_i$ is a single leaf. Thus
\begin{equation}\label{eq:sbound}
 0\le s\le \sum_{i=1}^k(m_i-1)=n-1-k.
\end{equation}
There is an important boundary distinction. If $k=n-1$, every $m_i=1$
and $s=0$. If $k\le n-2$, at least one $H_i$ is not a leaf, so $s\ge1$.
This explains both the missing value $k^2$ in a non-extremal slice and
the exceptional singleton at maximal depth.

\section{Every admissible pair is realized}
\label{sec:realization}
The bounds just obtained do not alone count the values: one must prove
that all intermediate $s$ occur. The following simple construction
makes that step explicit.

Let $R_r$ be the right-associated comb on $r$ leaves, defined by
$R_1=x$ and $R_{r+1}=x^{R_r}$. Let $C_r$ be the left-associated comb,
defined by $C_1=x$ and $C_{r+1}=(C_r)^x$.

\begin{lemma}[Prescribed left depth]\label[lemma]{lem:depth}
For every $m\ge2$ and $1\le h\le m-1$, an expression of size $m$
with left depth $h$ exists.
\end{lemma}
\begin{proof}
Start from $x^{R_{m-h}}$, and apply $h-1$ additional operations
$Y\mapsto Y^x$. The resulting expression has
\[
 1+(m-h)+(h-1)=m
\]
leaves and left depth $1+(h-1)=h$. All indicated indices are valid
because $m-h\ge1$.
\end{proof}

\begin{theorem}[Complete derivative-state classification]\label{thm:states}
Let
\[
 \Sset_n=\{(k(F),b(F)):F\in\E_n\}.
\]
Then $\Sset_1=\{(0,0)\}$, and for every $n\ge2$,
\begin{equation}\label{eq:states}
 \boxed{\quad
 \Sset_n=\{(n-1,(n-1)^2)\}
 \;\cup\!
 \bigcup_{k=1}^{n-2}
 \{(k,k^2+2s):1\le s\le n-k-1\}.
 \quad}
\end{equation}
In particular,
\begin{equation}\label{eq:actualset}
 \Dset_n=\{3(n-1)^2\}\cup
 \bigcup_{k=1}^{n-2}\{3k^2+6s:1\le s\le n-k-1\}
 \qquad(n\ge2).
\end{equation}
\end{theorem}
\begin{proof}
Necessity follows from \cref{lem:spine} and \eqref{eq:sbound}, together
with the distinction between $k=n-1$ and $k\le n-2$.

For sufficiency, the extremal pair is realized by $C_n$. Now fix
$1\le k\le n-2$ and $1\le s\le n-k-1$.
Choose $H_1$ of size $n-k$ and left depth $s$, using
\cref{lem:depth}, and put $H_2=\cdots=H_k=x$.
The expression in \eqref{eq:spine} then has
\[
 1+(n-k)+(k-1)=n
\]
leaves, left depth $k$, and
$\sum_i k(H_i)=s$. Its normalized third derivative is $k^2+2s$.
This realizes every pair on the right-hand side of \eqref{eq:states}.
\end{proof}

\begin{corollary}[Number of joint states]\label[corollary]{cor:joint}
For every $n\ge1$,
\begin{equation}\label{eq:jointcount}
 |\Sset_n|=1+\binom{n-1}{2}.
\end{equation}
\end{corollary}
\begin{proof}
For $n\ge2$, the different $k$-slices in \eqref{eq:states} have distinct
first coordinates, so their sizes add:
\[
 |\Sset_n|=1+\sum_{k=1}^{n-2}(n-k-1)
 =1+\frac{(n-1)(n-2)}2.
\]
For $n=1$ the same formula is $1$.
\end{proof}

\begin{remark}
Equation \eqref{eq:jointcount} does \emph{not} yet give $a(n)$.
Different second derivatives can accompany the same third derivative.
The first such collision occurs at $n=7$. The next section counts
these collisions exactly.
\end{remark}

\section{Counting the overlapping progressions}
\label{sec:overlaps}
For fixed $n$ and $1\le k\le n-2$, define
\begin{equation}\label{eq:Pk}
 P_k=\{k^2+2s:1\le s\le n-k-1\}.
\end{equation}
This is a step-$2$ integer interval with lower and upper endpoints
\begin{equation}\label{eq:endpoints}
 L_k=k^2+2,\qquad
 U_k=k^2+2n-2k-2=(k-1)^2+2n-3,
\end{equation}
and cardinality $|P_k|=n-k-1$.
The possible normalized third derivatives are their union together
with $\{(n-1)^2\}$.

The singleton is disjoint from all the progressions. For $n\ge3$,
$U_k$ is increasing in $k$, and
\[
 (n-1)^2-U_{n-2}=(n-1)^2-\bigl((n-2)^2+2\bigr)=2n-5>0.
\]
For $n=2$ there are no progressions.

Next, $P_k$ consists of integers congruent to $k$ modulo $2$.
Thus even-indexed and odd-indexed progressions are disjoint.
Within either parity, both endpoints strictly increase with $k$.

\begin{lemma}[Only the preceding same-parity interval matters]
\label[lemma]{lem:previous}
For $3\le k\le n-2$,
\begin{equation}\label{eq:intersection}
 \left|P_k\cap\bigcup_{j<k}P_j\right|
   =|P_k\cap P_{k-2}|=(n-3k+3)_+,
\end{equation}
where $(u)_+=\max(u,0)$.
\end{lemma}
\begin{proof}
Different parities cannot intersect. Among earlier progressions of the
same parity, $P_{k-2}$ has the largest upper endpoint.
If a value $v\in P_k$ lies in any earlier $P_j$ of the same parity,
then
\[
 L_{k-2}<L_k\le v\le U_j\le U_{k-2}.
\]
Its parity is that of $P_{k-2}$, so $v\in P_{k-2}$.
This proves the equality of the two intersections.

Since $U_{k-2}<U_k$ and $L_{k-2}<L_k$, a nonempty intersection runs
from $L_k$ to $U_{k-2}$ in steps of $2$. Its cardinality is
\begin{align*}
 \frac{U_{k-2}-L_k}{2}+1
 &=\frac{(k-2)^2+2n-2(k-2)-2-(k^2+2)}2+1\\
 &=n-3k+3.
\end{align*}
Taking the positive part also covers the empty case.
\end{proof}

This lemma is what makes triple and higher intersections harmless.
We add the progressions one at a time and subtract only those values
already present. A value occurring in several earlier progressions is
still subtracted just once when the next progression is inserted.
There is no assumption that pairwise overlaps are mutually disjoint.

\begin{proposition}[Exact overlap formula]\label[proposition]{prop:count}
For every $n\ge1$,
\begin{equation}\label{eq:overlapcount}
 \boxed{\qquad
 a(n)=1+\binom{n-1}{2}-\sum_{j\ge2}(n-3j)_+.
 \qquad}
\end{equation}
The sum is finite: only $j\le\lfloor(n-1)/3\rfloor$ contribute.
\end{proposition}
\begin{proof}
For $n\ge3$, sum the sizes of the progressions, subtract the previous
union intersection from \cref{lem:previous} each time, and then add
the disjoint singleton:
\[
 a(n)=1+\sum_{k=1}^{n-2}(n-k-1)
       -\sum_{k=3}^{n-2}(n-3k+3)_+.
\]
In the last sum put $j=k-1$. All subsequent terms are zero, so the
upper limit can be replaced by infinity. The first sum is
$\binom{n-1}{2}$. For $n=1,2$, the right-hand side of
\eqref{eq:overlapcount} is $1$, agreeing with the direct calculation.
\end{proof}

The factor $3$ in the overlap cutoff is structural. For consecutive
same-parity slices, the endpoint separation contains $3k-n$.
This, rather than any empirical fit, is the origin of the eventual
period $3$.

\section{Deriving the conjectured generating function}
\label{sec:gf}
We now sum \eqref{eq:overlapcount}. All calculations can be interpreted
in the ring of formal power series. In particular, exchanging the two
sums below is justified coefficient by coefficient: at a fixed degree,
only finitely many terms are nonzero.

The first two contributions are
\begin{equation}\label{eq:jointgf}
 \sum_{n\ge1}\left(1+\binom{n-1}{2}\right)z^n
 =\frac{z}{1-z}+\frac{z^3}{(1-z)^3}.
\end{equation}
For the correction, fix $j\ge2$. Substituting $n=3j+r$ gives
\[
 \sum_{n\ge1}(n-3j)_+z^n
   =z^{3j}\sum_{r\ge1}rz^r
   =\frac{z^{3j+1}}{(1-z)^2}.
\]
Therefore
\begin{equation}\label{eq:collisiongf}
 \sum_{n\ge1}\sum_{j\ge2}(n-3j)_+z^n
 =\frac{1}{(1-z)^2}\sum_{j\ge2}z^{3j+1}
 =\frac{z^7}{(1-z)^2(1-z^3)}.
\end{equation}
Subtracting \eqref{eq:collisiongf} from \eqref{eq:jointgf} proves
\eqref{eq:gfdecomp}. Notice that the correction begins at $z^7$,
exactly where collisions between different depths first appear.

For completeness, the algebraic conversion to the OEIS form is
\begin{align*}
 \GF
 &=\frac{(z-2z^2+2z^3)(1+z+z^2)-z^7}
 {(1-z)^3(1+z+z^2)}\\
 &=\frac{z-z^2+z^3+2z^5-z^7}
 {1-2z+z^2-z^3+2z^4-z^5}\\
 &=\frac{z(z^6-2z^4-z^2+z-1)}
 {z^5-2z^4+z^3-z^2+2z-1}.
\end{align*}
The final line multiplies numerator and denominator by $-1$.
Thus the stated rational generating function follows directly from
an exact count of all defining derivative values.

\section{The closed formula and its three residue classes}
\label{sec:closed}
We also evaluate \eqref{eq:overlapcount} explicitly. For $n\ge4$, put
\[
 q=\left\lfloor\frac{n-1}{3}\right\rfloor\ge1.
\]
Then
\[
 \sum_{j\ge2}(n-3j)_+
 =\sum_{j=2}^{q}(n-3j)
 =(q-1)n-3\left(\frac{q(q+1)}2-1\right).
\]
This identity includes the empty sum at $q=1$.
Substitution gives
\begin{equation}\label{eq:qformula}
 a(n)=\frac{n^2-n-2qn+3q^2+3q-2}{2}.
\end{equation}
Writing $n$ by its residue modulo $3$ yields
\begin{equation}\label{eq:residues}
 \begin{aligned}
 a(3m)&=3m^2-1, &&m\ge1,\\
 a(3m+1)&=3m^2+2m-1, &&m\ge1,\\
 a(3m+2)&=3m^2+4m, &&m\ge1.
 \end{aligned}
\end{equation}
For the first line at $m=1$, $a(3)=2$ is checked directly;
all other cases follow from \eqref{eq:qformula}.
Because the only square residues modulo $3$ are $0$ and $1$,
\eqref{eq:residues} is precisely
\[
 a(n)=\left\lfloor\frac{n^2-2}{3}\right\rfloor\qquad(n\ge3).
\]
Together with $a(1)=a(2)=1$, this finishes the proof of
\cref{thm:main}.

\subsection{Why the first two terms must be separated}
Substituting $n=1$ into the floor expression gives $-1$, and
substituting $n=2$ gives $0$. Neither is the desired value $1$.
A uniform version valid for every positive $n$ is
\begin{equation}\label{eq:delta}
 a(n)=\left\lfloor\frac{n^2-2}{3}\right\rfloor
       +2\,\mathbf1_{\{n=1\}}+\mathbf1_{\{n=2\}}.
\end{equation}
Equivalently,
\begin{equation}\label{eq:periodic}
 a(n)=\frac{n^2}{3}-\frac43+\frac13\mathbf1_{\{3\mid n\}}
       +2\,\mathbf1_{\{n=1\}}+\mathbf1_{\{n=2\}}.
\end{equation}
Thus a second, coefficient-based derivation of the generating function is
\begin{equation}\label{eq:gfindicator}
 \GF=2z+z^2+\frac{z(1+z)}{3(1-z)^3}
       -\frac{4z}{3(1-z)}+\frac{z^3}{3(1-z^3)}.
\end{equation}
Here the finite correction $2z+z^2$ accounts exactly for the two
exceptional initial terms. Simplifying \eqref{eq:gfindicator} gives
the same rational function as \eqref{eq:gfoeis}.

\section{Worked examples and the first cross-depth collision}
\label{sec:examples}
\subsection{Five copies of \texorpdfstring{$x$}{x}}
For $n=5$, the classification is small enough to display completely.
\begin{center}
\begin{tabular}{@{}cccc@{}}
\toprule
Left depth $k$ & Allowed $s$ & $b=F'''(1)/3$ & $F'''(1)$\\
\midrule
$1$ & $1,2,3$ & $3,5,7$ & $9,15,21$\\
$2$ & $1,2$ & $6,8$ & $18,24$\\
$3$ & $1$ & $11$ & $33$\\
$4$ & $0$ & $16$ & $48$\\
\bottomrule
\end{tabular}
\end{center}
The resulting set is
\[
 \Dset_5=\{9,15,18,21,24,33,48\},
\]
so $a(5)=7$. Within each row the second derivative is $2k$.
The accompanying file \texttt{witnesses\_n5\_n7.json} gives an explicit
fully parenthesized expression for every displayed joint state.

\subsection{Seven copies of \texorpdfstring{$x$}{x}}
At $n=7$, the normalized progressions and singleton are
\begin{center}
\begin{tabular}{@{}cl@{}}
\toprule
$k$ & Normalized third derivatives\\
\midrule
$1$ & $3,5,7,9,11$\\
$2$ & $6,8,10,12$\\
$3$ & $11,13,15$\\
$4$ & $18,20$\\
$5$ & $27$\\
$6$ & $36$\\
\bottomrule
\end{tabular}
\end{center}
There are $16$ joint states, but the value $b=11$ appears at both
$k=1$ and $k=3$. Hence there are only $15$ distinct third derivatives.
The two expressions
\[
 F_1=x^{C_6},\qquad F_2=\bigl((x^{R_4})^x\bigr)^x
\]
make the collision explicit. For $F_1$, $(k,s)=(1,5)$, while for
$F_2$, $(k,s)=(3,1)$. Both have $b=k^2+2s=11$, so
\[
 F_1'''(1)=F_2'''(1)=33,
 \qquad F_1''(1)=2,\quad F_2''(1)=6.
\]
They are therefore not the same analytic function, despite having the
same third derivative at $1$.

For $n\le6$, every term in the overlap correction is zero. Thus $n=7$
is the first collision \emph{between distinct joint derivative states}.
Different parenthesizations can already have identical joint states
at smaller sizes; these are a different kind of duplication.

\subsection{Initial values and a counting audit}
\begin{center}
\begin{tabular}{@{}rrrr@{}}
\toprule
$n$ & Joint states $1+\binom{n-1}{2}$ & Overlap correction & $a(n)$\\
\midrule
1&1&0&1\\
2&1&0&1\\
3&2&0&2\\
4&4&0&4\\
5&7&0&7\\
6&11&0&11\\
7&16&1&15\\
8&22&2&20\\
9&29&3&26\\
10&37&5&32\\
11&46&7&39\\
12&56&9&47\\
13&67&12&55\\
14&79&15&64\\
15&92&18&74\\
\bottomrule
\end{tabular}
\end{center}
The archive includes this audit through $n=60$ and the sequence
coefficients through $n=1000$.

\section{Consequences of the classification}
\label{sec:consequences}
\subsection{Extremal derivatives}
\begin{corollary}\label[corollary]{cor:extrema}
For $n\ge3$,
\[
 \min\Dset_n=9,\qquad \max\Dset_n=3(n-1)^2.
\]
The maximum is attained only by the fully left-associated expression
$C_n$. For $n=1,2$, the derivative sets are $\{0\}$ and $\{3\}$,
respectively.
\end{corollary}
\begin{proof}
In a non-extremal slice, $k\ge1$ and $s\ge1$, so
$k^2+2s\ge3$, with equality attained at $k=s=1$.
The exceptional singleton is larger than every other normalized
value, as shown in \cref{sec:overlaps}. Attaining it forces
$k=n-1$, so every right subtree along the spine is a leaf.
That is exactly $C_n$.
\end{proof}

For $n\ge3$, the minimum can be described just as explicitly:
it is attained exactly by expressions $x^{(x^H)}$ with
$H\in\E_{n-2}$. Indeed, $k=1$ means $F=x^G$, and $s=1$ means that
$G$ has left depth $1$, hence $G=x^H$.

\subsection{Simple recurrences}
From the closed formula, for $n\ge3$,
\begin{equation}\label{eq:step3}
 a(n+3)=a(n)+2n+3.
\end{equation}
This follows because adding $3$ to $n$ adds the integer $2n+3$
to $(n^2-2)/3$, so it adds the same integer to its floor.
The forward differences are
\begin{equation}\label{eq:step1}
 a(n+1)-a(n)=\left\lceil\frac{2n}{3}\right\rceil
 \qquad(n\ge3),
\end{equation}
as is checked directly in the three residue classes in
\eqref{eq:residues}.
The stated starting indices matter: these recurrences do not extend
unchanged across both exceptional initial terms.

Expanding the denominator in \eqref{eq:gfpositive} also gives the
homogeneous recurrence
\begin{equation}\label{eq:order5}
 \begin{split}
 a(n)={}&2a(n-1)-a(n-2)+a(n-3)\\
        &-2a(n-4)+a(n-5),\qquad n\ge8.
 \end{split}
\end{equation}
The initial data are
$a(1),\ldots,a(7)=1,1,2,4,7,11,15$.
The numerator has degree $7$, which explains the onset $n=8$.

\subsection{Quadratic growth and exact periodic correction}
For $n\ge3$, \eqref{eq:periodic} gives the exact identity
\begin{equation}\label{eq:asymptotic}
 a(n)=\frac{n^2}{3}-\frac43+\frac13\mathbf1_{\{3\mid n\}}.
\end{equation}
In particular,
\[
 a(n)=\frac{n^2}{3}+O(1),\qquad a(n)\sim\frac{n^2}{3}.
\]
The constant correction is $-1$ on multiples of $3$ and $-4/3$
on the other two residue classes. Thus the sequence is eventually
a quadratic quasipolynomial of period $3$; it is not eventually a
single polynomial.

\subsection{Why the denominator has degree five}
The denominator factors as
\[
 (1-z)^2(1-z^3)=(1-z)^3(1+z+z^2).
\]
The cubic factor at $z=1$ corresponds to quadratic growth. The two
primitive cube roots of unity encode the periodic correction.
These factors do not cancel. For
\[
 N(z)=z(1-z+z^2+2z^4-z^6),
\]
we have $N(1)=2$. If $\omega^3=1$ and $\omega\ne1$, then
$N(\omega)=-\omega\ne0$.

Consequently, \eqref{eq:order5} has the smallest possible order among
eventual homogeneous recurrences with constant coefficients.
Indeed, any such recurrence would give a polynomial $Q(z)$ of the
same order for which $Q(z)\GF$ is a polynomial. The reduced
denominator of $\GF$ must divide $Q(z)$, so $\deg Q\ge5$.
The shorter relation \eqref{eq:step3} is inhomogeneous and therefore
does not contradict this statement.

\section{Exact computation and reproducibility}
\label{sec:verification}
The proof above is valid for every $n$ and does not rely on finite
verification. The archive nevertheless supplies independent exact
checks, both to make the result reproducible and to catch implementation
or indexing errors.

\subsection{Three independent routes to the data}
The companion \texttt{verify.py} uses only Python's standard library.
It implements the following mathematically distinct computations.

\paragraph{Recursive derivative-state enumeration.}
Starting from $\{(0,0)\}$ at size $1$, it combines all splits
$n=j+(n-j)$ using \eqref{eq:state-law}. This computation does not use
the claimed explicit state set or count. Because the right child's
$b$-coordinate is irrelevant, its actually computed depth projection
is used to avoid redundant work.

\paragraph{Exhaustive expression enumeration with Taylor arithmetic.}
All ordered binary parenthesizations are enumerated up to a selected
small size. Each is evaluated in $\mathbb Q[t]/(t^4)$ after replacing
$x$ by $1+t$, using
\[
 G^H=\exp(H\log G),\qquad
 \log(1+w)=w-\frac{w^2}{2}+\frac{w^3}{3},\qquad
 e^w=1+w+\frac{w^2}{2}+\frac{w^3}{6}
 \pmod{t^4}.
\]
This evaluator does not use the derivative-state recurrence.
If the resulting jet is
$1+t+c_2t^2+c_3t^3$, then $(k,b)=(c_2,2c_3)$.
The coefficients are exact rational numbers, not finite-difference
approximations.

\paragraph{Expansion of the proposed rational function.}
The script separately expands the original OEIS numerator and
denominator by coefficient division. Its coefficients are checked
against both the floor formula and the overlap formula. This route
is also used to check \eqref{eq:step3} and \eqref{eq:order5}.

The script additionally constructs a witness for every admissible
joint state up to a selected size, counts its leaves, and checks its
Taylor jet with the independent evaluator.

\subsection{Executed checks}
The supplied \texttt{verification\_report.json} records a successful run
with the following bounds:
\begin{center}
\begin{tabular}{@{}p{0.63\linewidth}r@{}}
\toprule
Check & Executed bound\\
\midrule
Recursive states versus the complete classification & $1\le n\le60$\\
All parenthesizations, evaluated by exact Taylor jets & $1\le n\le10$\\
Constructive witnesses, each checked by Taylor jets & $1\le n\le20$\\
Rational g.f. coefficients versus both count formulas & $1\le n\le1000$\\
\bottomrule
\end{tabular}
\end{center}
There were $4{,}862$ parenthesizations at $n=10$ and $6{,}918$
in total through $n=10$. All $1{,}160$ constructed witnesses through
$n=20$ passed their leaf-count and derivative checks.
No floating-point differentiation or numerical tolerance was used.

For context, if $c_n=|\E_n|$, the recursive syntax gives
$c_1=1$ and $c_n=\sum_{j=1}^{n-1}c_jc_{n-j}$.
Hence $C(z)=\sum_{n\ge1}c_nz^n$ satisfies $C=z+C^2$, and so
$c_n=\binom{2n-2}{n-1}/n$.
The exponential number of expressions explains why direct symbolic
enumeration is useful for small independent checks but not needed
for the proof or for large $n$.

\subsection{How to reproduce the files}
From the archive directory, run
\begin{lstlisting}[language=bash]
python3 verify.py --out-dir artifacts
pdflatex -interaction=nonstopmode -halt-on-error article.tex
pdflatex -interaction=nonstopmode -halt-on-error article.tex
\end{lstlisting}
The Python command has no third-party dependencies. Use Python 3.10 or
later, without the \texttt{-O} option, since verification uses assertions.
The PDF is built with a standard TeX Live or MiKTeX installation that
includes the packages declared at the top of \texttt{article.tex}.
A \texttt{Makefile} and detailed \texttt{README.md} are also supplied.

For a single $n$, the full state set is available by the short formula
\begin{lstlisting}[language=Python]
states = {(n - 1, (n - 1)**2)} | {
    (k, k*k + 2*s)
    for k in range(1, n - 1)
    for s in range(1, n - k)
}
third_derivatives = sorted({3*b for k, b in states})
\end{lstlisting}
This formula also handles $n=1$ and $n=2$ correctly. It uses
$O(n^2)$ state insertions for one size. By contrast, evaluating the
count itself needs only a constant number of integer arithmetic
operations via \eqref{eq:closed}; this is not a claim of constant
bit complexity for unbounded integers.

\section{Scope of the result}
The original task is resolved completely: \eqref{eq:gfoeis} is proved
for the sequence as defined, including its offset and initial terms.
The stronger classification \eqref{eq:states} supplies explicit
witnesses and rules out missing derivative values. The argument counts
distinct values, not their frequencies among parenthesizations.

Higher derivative counts form the broader table recorded as OEIS
A216368 \cite{OEIS216368}. The present proof makes no assertion that
the fourth- or higher-derivative sequences have analogous rational
generating functions. The fact that only two integer statistics are
needed here is a special feature of the third-order calculation.

\appendix
\section{A second derivation of the spine formula}
\label{app:alternate}
There is a useful analytic check on \cref{lem:spine} that bypasses
successive use of the derivative recurrence. On $(0,\infty)$, the
ordinary identity $(u^v)^w=u^{vw}$ allows \eqref{eq:spine} to be written
\[
 F(x)=x^{P(x)},\qquad P(x)=\prod_{i=1}^kH_i(x).
\]
At $x=1$, each $H_i$ and $H_i'$ equals $1$. With
$s=\sum_i k(H_i)$, the product rule gives
\[
 P(1)=1,\qquad P'(1)=k,\qquad
 P''(1)=\sum_iH_i''(1)+\sum_{i\ne j}H_i'(1)H_j'(1)
       =2s+k(k-1).
\]
Set $V=P\log x$. Then
\begin{align*}
 V'(1)&=1,\\
 V''(1)&=2k-1,\\
 V'''(1)&=3P''(1)-3P'(1)+2
          =6s+3k^2-6k+2.
\end{align*}
Since $F=e^V$, we obtain
\[
 F''(1)=V''(1)+(V'(1))^2=2k
\]
and
\[
 F'''(1)=V'''(1)+3V'(1)V''(1)+(V'(1))^3=3k^2+6s.
\]
The cancellation of the linear and constant terms makes the square
$k^2$ especially transparent in this form.

\clearpage
\section{A compact proof suitable for a sequence note}
\label{app:compact}
For each parenthesized expression $F$, define
$k(F)=F''(1)/2$ and $b(F)=F'''(1)/3$.
Direct differentiation gives
\[
 k(G^H)=k(G)+1,\qquad
 b(G^H)=b(G)+2k(G)+2k(H)+1,
\]
with $k(x)=b(x)=0$.
Along a left spine of length $k$ with right subexpressions $H_i$,
this gives $b=k^2+2\sum_i k(H_i)$.
For size $n$, the resulting joint states are exactly
\[
 (n-1,(n-1)^2),\qquad
 (k,k^2+2s)\quad(1\le k\le n-2,\ 1\le s\le n-k-1).
\]
For realizability, put all $n-k$ nontrivial right-subtree leaves into
one subtree of depth $s$ and make the remaining right subtrees leaves;
a size-$m$ tree of any left depth $1,\ldots,m-1$ is obtained by
successively adding left-spine nodes to a right comb.

At fixed $k$, the possible $b$ form a step-$2$ interval with endpoints
$k^2+2$ and $(k-1)^2+2n-3$. Only equal parities overlap. On inserting
the $k$th interval, its intersection with the preceding union is
exactly its intersection with the $(k-2)$nd interval, of size
$(n-3k+3)_+$. The extremal singleton is disjoint. Therefore
\[
 a(n)=1+\binom{n-1}{2}-\sum_{j\ge2}(n-3j)_+.
\]
Summing over $n\ge1$ gives
\[
 \sum_{n\ge1}a(n)z^n
 =\frac{z}{1-z}+\frac{z^3}{(1-z)^3}
  -\frac{z^7}{(1-z)^2(1-z^3)},
\]
which simplifies to the conjectured rational function.
Evaluating the finite sum gives
$a(n)=\lfloor(n^2-2)/3\rfloor$ for $n\ge3$, with $a(1)=a(2)=1$.

\begin{thebibliography}{9}
\bibitem{OEIS199085}
The OEIS Foundation Inc.,
\emph{A199085: Number of distinct values taken by the third derivative
of parenthesized self-exponentials at $x=1$},
\url{https://oeis.org/A199085}.
Text endpoint supplied with the problem:
\url{https://oeis.org/search?q=id:A199085&fmt=text}.
Accessed September 20, 2026.

\bibitem{OEIS216368}
The OEIS Foundation Inc.,
\emph{A216368: Counts of distinct derivative values of parenthesized
self-exponentials},
\url{https://oeis.org/A216368}.
Accessed September 20, 2026.
\end{thebibliography}
\end{document}
